\documentclass[12pt, a4paper]{academics}

\academicsCourse{数据库原理与应用}{Database Principle and Application}
\academicsReport{概念结构设计——ER 模型图}{2026 年 4 月 2 日}
\academicsArthur{华博文}{19240212}

\begin{document}

\makeacademicscover

\section{实验环境}

\subsection{操作系统}

macOS 26.4 arm64，Darwin 25.4.0，Apple M5。

\subsection{绘图工具}

sql\_to\_ER（GitHub: systemsrx/sql\_to\_ER），从 SQL DDL 自动生成 ER 模型图。

\subsection{文档工具}

\LaTeX{}。

\section{实验目的}

\begin{enumerate}
    \item 理解概念结构设计在数据库设计全流程中的核心地位——它是从需求分析到逻辑结构设计的桥梁，将模糊的业务描述转化为精确的数据模型；
    \item 掌握从需求说明中提取实体、识别实体间联系（1:1、1:N、M:N）以及区分弱实体与关联实体的方法；
    \item 能够使用 ER 模型图规范地表达概念数据模型，为后续逻辑结构设计与物理实现奠定基础。
\end{enumerate}

\section{实验内容}

\subsection{需求概述}

本次实验以 B2C 在线购物平台为背景。平台服务日均活跃用户约 10 万人，商品总量约 10 万，日均订单量约 1 万笔。围绕「浏览—挑选—下单—收货」这一核心购物流程，系统需支持以下功能：

\begin{itemize}
    \item \textbf{用户管理}：顾客注册账户，提供用户名、密码、真实姓名、联系电话、收货地址等信息，可管理多个收货地址。
    \item \textbf{商品管理}：商品包含唯一编号、名称、描述、价格、类别、库存数量等。
    \item \textbf{购物车功能}：顾客可将商品加入购物车并指定数量；购物车数据在登录状态下持久化；数量不能超库存。
    \item \textbf{订单处理}：从购物车结算生成订单，含订单号、下单时间、总金额、收货地址、支付状态、订单状态等信息；订单提交时扣减库存，取消未发货订单时恢复库存。
    \item \textbf{支付与物流}：模拟支付流程；顾客可查看订单进度及物流信息。
    \item \textbf{查询与统计}：支持按关键词、分类、价格区间搜索商品，按销量、价格、上架时间排序。
\end{itemize}

\subsection{实体名称列表}

对上述需求进行分析，沿购物流程的四个环节逐层识别，共提取出 7 个实体。其中订单明细为弱实体（依赖于订单而存在），购物车项为关联实体（解析顾客与商品之间的 M:N 联系）。物流状态（待发货/已发货/已完成/已取消）作为订单的属性处理，不单独建模。

\begin{table}[h]
\centering
\caption{实体名称列表}
\setlength{\belowcaptionskip}{8pt}
\begin{tabular}{|c|c|c|p{6cm}|}
\hline
\textbf{序号} & \textbf{实体名称} & \textbf{类型} & \textbf{说明} \\
\hline
1 & 顾客（Customer） & 核心实体 & 平台注册的消费者，拥有登录凭据、个人信息及多个收货地址 \\
2 & 类别（Category） & 核心实体 & 商品所属的分类，支持按分类检索与浏览 \\
3 & 商品（Product） & 核心实体 & 平台待售商品，记录名称、描述、价格、库存等信息 \\
4 & 购物车项（CartItem） & 关联实体 & 顾客与商品之间 M:N 联系的关联实体，携带购买数量 \\
5 & 收货地址（Address） & 核心实体 & 顾客的收货地址，一个顾客可拥有多个地址，订单引用一个地址 \\
6 & 订单（Order） & 核心实体 & 顾客提交的购买订单，记录下单时间、总金额、支付与物流状态等 \\
7 & 订单明细（OrderItem） & 弱实体 & 订单的明细行，记录每项商品的快照（名称、单价、数量）；主键包含订单ID，离开所属订单无独立标识 \\
\hline
\end{tabular}
\end{table}

\subsection{ER 模型图}

采用 Chen 表示法绘制 ER 模型图（不绘制实体属性）。矩形表示实体，双线矩形表示弱实体，菱形表示联系，连线上标注基数约束。

\begin{figure}[h]
\centering
\includegraphics[width=\textwidth]{../res/er-diagram.png}
\caption{B2C 在线购物平台 ER 模型图（Chen 表示法）}
\end{figure}

\subsection{实体间关系说明}

各实体间的联系及基数约束如下：

\begin{enumerate}
    \item \textbf{类别——商品（1:N）}：一个类别可包含多个商品，一个商品仅属于一个类别。这是商品浏览的基础——顾客通过类别逐层定位目标商品。
    \item \textbf{顾客——地址（1:N）}：一个顾客可管理多个收货地址（如家庭地址、公司地址），一个地址仅属于一个顾客。下单时可从中选择一个作为本次收货地址。
    \item \textbf{顾客——商品，通过购物车项（M:N）}：需求中「顾客将商品加入购物车」描述了一个典型的 M:N 场景。购物车项将这一联系拆解为商品 1:N 购物车项和购物车项 N:1 顾客两个 1:N 联系，并携带购买数量属性，承担了「数量不得超过当前库存」的约束校验责任。
    \item \textbf{顾客——订单（1:N）}：一个顾客可下多个订单，一个订单仅属于一个顾客。订单记录了顾客在某时刻的购买快照——即使顾客后续修改了个人信息或地址，订单中的收货地址和历史价格不受影响。
    \item \textbf{订单——地址（N:1）}：每笔订单引用一个收货地址，同一个地址可被多笔订单使用。这一联系将订单与配送目的地绑定，且独立于「顾客拥有地址」的联系——即便地址被删除，订单中引用的地址快照仍应保留。
    \item \textbf{订单——订单明细（1:N，弱实体）}：订单由若干明细项构成。OrderItem 是 Order 的弱实体，其存在依赖于所属订单。订单提交时，订单明细中的商品名称、单价等字段为下单时刻的快照值，而非对 Product 表的实时引用——这是为了在商品信息变更后仍能追溯历史订单的准确内容。
    \item \textbf{商品——订单明细（1:N）}：一个商品可出现在多个订单的明细中，但一条订单明细仅对应一种商品。此联系保证了从商品维度追溯销售记录的能力。
\end{enumerate}

\subsection{关键设计决策}

在设计 ER 模型的过程中，以下四个决策点值得说明：

\textbf{物流状态为何不单独建模？} 需求中提到了「物流信息」和订单状态（待发货/已发货/已完成），两者的粒度高度一致——物流状态的变更本质上就是订单状态的推进。若将物流建模为物流单弱实体，反而引入了不必要的概念复杂性。将物流状态作为订单的属性（如订单状态、运单号、承运商）处理更加简洁，且满足当前需求的追溯粒度。只有在需要记录每一条物流轨迹（如分拨中心扫描记录）时，物流单独立建模才有必要，而本次需求未涉及此粒度。

\textbf{收货地址为何作为独立实体而非顾客的属性？} 需求明确要求顾客可「管理多个收货地址」，且订单需引用具体的收货地址。若将其设计为顾客的多值属性，则订单无法直接引用某条地址——这正是将多值属性提升为独立实体的经典场景。

\textbf{订单明细为何采用快照而非引用？} 订单明细记录了「商品名称、单价、数量」等字段，这些是下单时刻的冗余副本，而非对 Product 表的纯外键引用。这一设计决策源于以下考虑：商品的价格和名称可能在订单存续期内发生变化，若订单明细仅存储 product\_id 后再关联查询，则历史订单将显示修改后的价格——这在支付、退款、审计等场景中不可接受。

\textbf{购物车为何建模为关联实体？} 购物车项连接顾客与商品两个实体，并携带购买数量这一联系属性——该属性既不属于顾客也不属于商品，只能挂靠在联系本身。此外，购物车项同时与顾客和商品保持独立的 1:N 联系，具备独立标识（购物车项ID），因此将其建模为关联实体而非多值属性，符合 ER 建模规范中对 M:N 联系的标准处理方式。

\section{实验过程归纳总结}

本次实验的核心任务是从自然语言需求描述中识别出 7 个实体及 7 组联系。回顾这一过程，有以下三点方法论层面的收获。

\textbf{第一，业务流程驱动的实体识别。}面对一段完整的自然语言需求，最有效的切入方式不是通读后凭直觉列实体，而是沿业务流程的主线——浏览、挑选、下单、收货——逐环节追问「这个环节涉及了什么？谁做的？产生了什么记录？」。这一逐环节分解的策略能够最大程度地避免遗漏。例如，「收货地址」在初读时容易被忽略（它出现在「用户管理」部分，远离「下单」环节），但沿流程追问「下单时收货地址从何而来」便能发现它具备独立标识和独立生命周期，因此必须建模为独立实体。

\textbf{第二，M:N 联系的关联实体化。}需求中「顾客将商品加入购物车」描述了一个典型的 M:N 场景。在概念模型层面，这种联系不能简单地画一条两端标注 M 和 N 的连线了事——因为购买数量这一属性无法挂在连线本身，必须引入关联实体来承载。关联实体同时将 M:N 联系拆分为两个 1:N 联系，使概念模型能够无损映射到关系模式。这一转化是 ER 建模中最常见也最容易被忽视的环节。

\textbf{第三，建模粒度的权衡——少即是多。}物流状态最初被考虑为独立实体，但分析发现其信息粒度与订单状态高度同构——物流状态的推进本质上就是订单状态的变更。这种情况下强行拆分为独立实体不仅增加模型复杂度，还会在逻辑设计阶段引入不必要的关联表。此次实验最深刻的体会是：不是需求中出现的每个名词都值得成为实体，只有那些具备独立标识和独立生命周期的概念才应当独立建模。这一判断力是概念结构设计中最需要培养的能力。

\section{结论}

本实验完成了 B2C 在线购物平台数据库的概念结构设计，通过对需求描述的逐层分析，识别出顾客、商品、订单等 7 个核心实体，以及包含(1:N)、拥有(1:N)、下单(1:N)、收货地址(N:1)等 7 组联系，绘制了规范的 ER 模型图（Chen 表示法）。模型完整覆盖了「浏览—挑选—下单—收货」的核心购物流程，在实体粒度上做到了适度抽象——物流信息作为订单属性而非独立实体——并在收货地址独立建模、订单明细快照策略、购物车关联实体设计等关键决策点上给出了合理的设计依据，为后续的逻辑结构设计提供了清晰的概念基础。

\appendix
\section{附录：数据库模式 SQL}

以下为 ER 模型对应的关系模式 DDL，共 7 张表，与实体列表一一对应。

\inputminted{text}{../res/schema.sql}

\end{document}
